\title{Byte Latent Transformer: Patches Scale Better Than Tokens}

\begin{abstract}
We present the Byte Latent Transformer ({\textsc BLT}), a novel large language model architecture operating at the byte level, which for the first time attains performance comparable to that of traditional tokenization-based LLMs at scale, while substantially enhancing inference speed and resilience. 
{\textsc BLT} represents input as patches of bytes whose sizes are determined dynamically; these patches become the core computational elements of the model. Patch boundaries are informed by the entropy associated with predicting the next byte, so computational resources and representational power are focused on regions with greater data intricacy. We conduct the first scaling analysis of byte-level models, rigorously controlling floating point operations (flops), extending up to 8B parameters and training on 4T raw bytes. Our findings establish that large-scale models can be trained directly on raw byte sequences, dispensing with pre-defined vocabularies. Both the learning and inference processes are optimized via adaptive selection of longer patches in low-entropy, more predictable contexts, leading to improved performance on tasks requiring advanced reasoning and handling infrequent cases. Notably, at equivalent inference cost, {\textsc BLT} exhibits superior scaling relative to tokenization-based approaches, achieved through the joint expansion of both patch size and model capacity.
\end{abstract}

\section{Introduction}
\label{section:intro}

We present the Byte Latent Transformer~(\textbf{BLT}), a model that operates directly on raw byte sequences without relying on tokenization. BLT represents the first instance where a tokenizer-free approach achieves comparable results to traditional token-based architectures at scale, while also delivering notable gains in efficiency and resilience (\S\ref{sec:robustness}).

While current large language models (llms) are predominantly trained in an end-to-end manner, they still depend on tokenization—a pre-processing step employing heuristics to segment byte sequences into a predetermined token inventory. This fixed token representation can distort string encoding, resulting in issues such as increased sensitivity to specific domains and modalities~\citep{dagangetting}, vulnerability to input perturbations~(\S\ref{sec:robustness}), insufficient modeling of orthographic features~\citep{edman2024cute}, and unequal treatment across languages~\citep{liang2023xlm,petrov2024language,limisiewicz2024myte}.

Historically, tokenization has played a crucial role since training llms directly on raw bytes incurs significant computational expense due to the resulting extended sequence lengths~\citep{xue2022byt5}. Previous research has attempted to address this by introducing self-attention mechanisms with improved efficiency~\citep{el2020characterbert, clark2022canine} or by adopting architectures that do not rely on attention~\citep{wang2024mambabyte}~(\S\ref{section:related_work}). Nonetheless, these strategies are mainly advantageous when developing \textit{small models}. When scaling up, the primary computational bottleneck in a Transformer arises from the extensive feed-forward network layers, which must be processed for every byte, rather than from the attention component itself.

For effective compute allocation, we introduce a dynamic and trainable approach that organizes bytes into \textit{patches}~(\S\ref{section:patching}), coupled with a novel architecture designed to integrate information from both individual bytes and patches. In contrast to tokenization, {\textsc BLT} does not employ a static patch vocabulary. Instead, any sequence of bytes can be transformed into a latent patch representation through compact, learnable encoder and decoder components. Our results demonstrate that this approach provides \textit{greater} computational efficiency compared to models reliant on tokenization.

Tokenization-based language models assign a uniform amount of computation to each token, resulting in a tradeoff between computational efficiency and predictive performance. This uniform allocation arises because tokens are derived using compression schemes that do not necessarily correspond to the complexity involved in making predictions. The foundation of our proposed architecture is the principle that computation should be allocated adaptively, responding to specific demands within the data. For instance, predicting the final characters of words usually presents lower complexity and entropy, suggesting that extensive transformer models are superfluous for such tasks—unlike the selection of the initial word in a sentence, which is typically more challenging. Accordingly, {\textsc BLT} features an architecture (\S\ref{section:architecture}) composed of three transformer blocks: two lightweight, byte-level \textit{local models} complemented by a larger, comprehensive global \textit{latent transformer}~(\autoref{fig:arch_high_level}). 

In {\textsc BLT}, the allocation of compute and formation of byte patches are governed by segmenting input based on the entropy associated with next-byte prediction. This approach yields contextualized byte groupings that maintain a relatively consistent information density.

We report the first scaling investigation of byte-level models with flop-based controls, scaling up to 8B parameters and 4T bytes of training data. This demonstrates that it is feasible to train large-scale models directly from bytes, entirely bypassing fixed-vocabulary tokenization. Overall, {\textsc BLT} achieves comparable training performance per flop\footnote{We measure computational cost by tallying the number of Floating Point OPerations (flops) expended by the model.} to Llama 3, yet achieves up to a 50\% reduction in inference flops~(\S\ref{section:scaling}). Furthermore, we show that processing inputs at the raw byte level yields notable gains for modeling rare or infrequent data patterns. In our tests, {\textsc BLT} models outperform tokenization-based alternatives in their robustness to noisy data and excel in tasks requiring fine-grained character-level reasoning, such as orthographic processing, phonological tasks, and machine translation for low-resource languages~(\S\ref{sec:robustness}). 

Additionally, {\textsc BLT} architectures make it possible to scale up both the model and patch size, holding inference flops constant. Increasing the patch size typically reduces computation, permitting more resources to be reallocated to expanding the global latent transformer component, as it is invoked less frequently. Our experiments with inference-flop controlled scaling~(\autoref{fig:fixed_inference_scaling}) reveal clear advantages in scaling characteristics compared to approaches reliant on tokenization.

To conclude, this work offers several key advancements: 1) We present {\textsc BLT}, a byte-level latent language model architecture capable of adaptive compute allocation to enhance flop efficiency; 2) We demonstrate that our approach matches the flop-controlled training performance of Llama 3 at scales up to 8B parameters, and can yield up to a 50\% improvement in flop efficiency by tolerating minimal reductions in evaluation metrics; 3) {\textsc BLT} models enable a new scaling regime for large language models, allowing growth in model size without exceeding a constant inference compute limit; 4) We provide evidence that {\textsc BLT} architectures exhibit greater resilience to noisy inputs and possess sub-word sensitivity that traditional token-based llms do not capture. The source code for training and inference with {\textsc BLT} is publicly available at~\url{https://github.com/facebookresearch/blt}.

\section{Patching: From Individual Bytes to Groups of Bytes}
\label{section:patching}

Dividing bytes into \textit{patches} enables {\textsc BLT} to flexibly assign computational resources according to the context. Figure~\ref{fig:patching_types} presents a range of strategies for partitioning bytes into patches. More formally, a patching function $f_p$ transforms an input byte sequence $\pmb{x} = \{x_i,|i=1,\ldots, n\}$ of length $n$ into a patch sequence $\pmb{p} = \{p_j|j=1,\ldots,m\}$ containing $m < n$ patches by mapping every $x_i$ onto \{0,1\}, where a value of 1 signals the initiation of a new patch. For both token-based and patch-based architectures, the predominant driver of computational demand is the number of main Transformer steps necessary to handle the input. In {\textsc BLT}, this corresponds to the total patches produced via a selected patching function. Thus, the key metric governing the data processing cost for training or inference with any patching method is the mean number of bytes per patch—referred to as the \textit{patch size}~(\S\ref{section:flops}). 

We proceed by describing three specific patching mechanisms: fixed-length byte patching~(\S\ref{section:static-patch}), whitespace-based patching~(\S\ref{section:white-patch}), and adaptive patching informed by local entropies from a compact byte language model~(\S\ref{section:dyn-patch}). In addition, we examine the process of incremental patching and highlight distinctions between patching and tokenization~(\S\ref{section:bpe-patch}).

\subsection{Strided Patching Every K Bytes}
\label{section:static-patch}

A commonly used approach for organizing bytes is to divide them into equally sized patches of $k$ bytes, as implemented in MegaByte~\citep{yu2023megabyte}. Using a fixed stride is both computationally simple for training and inference, and allows one to easily adjust the typical patch size, making control over computational requirements transparent. Nevertheless, this method introduces several notable limitations. Firstly, it lacks adaptive allocation of computational effort; for instance, computational resources might be squandered when a transformer step $j$ is dedicated solely to predicting uninformative elements like whitespace in code, while regions containing complex, information-dense content (such as mathematics) may receive insufficient attention. Secondly, it causes non-uniform, non-context-aware segmentation of byte sequences, so identical strings may be fragmented differently depending on their position.

\subsection{Space Patching}
\label{section:white-patch}

\citet{slagle2024spacebyte} introduces an enhanced variant of strided patching in which new patches are formed immediately following any space-like byte\footnote{
    Space-like bytes encompass all bytes except for Latin characters, digits, or utf-8 continuation bytes. Moreover, each patch must include at least one byte that is not space-like. }, leveraging the fact that such bytes commonly coincide with linguistic unit boundaries across languages. Under the Space patching scheme, a latent transformer operation (i.e., increased computational steps) is dedicated to every individual word, thereby maintaining consistent patching for words across different sequences. This approach also concentrates computational effort where predictions are most challenging—typically immediately after spaces. As an illustrative case, consider the task of generating the first byte in response to ``Who composed the Magic Flute?\uline{\hspace{.6em}}''. Choosing this initial character is markedly more complex due to a broader set of possibilities, whereas generating subsequent characters—following the "M" in ``Mozart''—is considerably easier. Nonetheless, Space patching displays limitations when confronted with certain languages or domains and lacks flexibility in adjusting patch sizes. To address these constraints, we proceed to present a novel patching method that both harnesses the observation that initial bytes of words present the most prediction difficulty and introduces a flexible framework for dynamically controlling patch sizes.

\subsection{Entropy Patching: Using Next-Byte Entropies from a Small Byte LM}
\label{section:dyn-patch}

In place of employing a rule-based strategy, for example using whitespace as a delimiter, we adopt a data-centric methodology to pinpoint regions with elevated uncertainty in next-byte predictions. Our method, termed \textit{entropy patching}, utilizes entropy measurements to establish the borders of each patch.

Specifically, we fit a compact byte-level auto-regressive language model using the {\textsc BLT} training data, and for each position, we determine the entropy of the upcoming byte under the LM’s probability distribution $p_{e}$ across the byte vocabulary $\mathcal{V}$:

\begin{equation}
H(x_i) = \sum_{v \in \mathcal{V}} p_{e}(x_i=v|\pmb{x}_{<i}) \log p_{e}(x_i=v|\pmb{x}_{<i})
\end{equation}

We explore two approaches for detecting patch boundaries using the entropies $H(x_i)$. The first approach selects locations where the entropy exceeds a fixed global threshold, as depicted in~\autoref{fig:patching}. The second approach marks points where the entropy is significantly larger compared to its immediate predecessor. This latter method may also be viewed as pinpointing instances where the local pattern of roughly monotonic entropy decrease within a patch is disrupted.

\begin{align*}
    \text{Global Constraint}\hspace{1cm}H(x_t)& > \theta_g\\
    \text{Approx. Monotonic Constraint}\hspace{1cm}H(x_t)& - H(x_{t-1}) > \theta_r
\end{align*}

Patch boundaries are determined in a minimal preprocessing stage carried out while loading the data. This contrasts with the approach of~\citet{nawrot-etal-2023-efficient}, who utilize a classifier specifically trained to detect entropy-derived patch boundaries. In our experiments~(\S\ref{section:experiments}), we evaluate and contrast these two strategies for discriminating between bytes of low and high entropy.

\subsection{The Byte-Pair Encoding (BPE) Tokenizer and Incremental Patching}
\label{section:incremental-patching}
\label{section:bpe-patch}

Numerous contemporary llms, such as our reference Llama 3, employ subword tokenizers such as bpe~\citep{gage1994new, sennrich-etal-2016-neural}. Throughout this work, we denote ``tokens'' as byte sequences selected from a \textit{finite} set established before training, distinguishing them from ``patches,'' which are adaptively formed sequences lacking a predetermined vocabulary. An essential distinction is that, when using tokens, the model cannot directly interact with the raw byte features, whereas patches do not impose this limitation.

An important advancement introduced by {\textsc BLT} compared to tokenization-based architectures is its reconfiguration of the balance between vocabulary size and computational requirements. Traditional large language models must accept that a bigger vocabulary leads to tokens that encode more information on average, thereby reducing the total number of model steps during inference. However, this comes at the expense of increasing the output dimensionality of the model’s final projection layer, which can greatly raise computational demands. As a result, tokenization-dependent methods are constrained in their ability to flexibly adjust token length or the corresponding inference costs. To illustrate, Llama 3 boosts the mean token size from 3.7 to 4.4 bytes, which is accomplished by expanding the size of its embedding matrix by a factor of four relative to Llama 2.

During the generation process, {\textsc BLT} must determine at each position in the byte sequence whether it is at the start of a new patch. This decision dictates if additional computation involving the Latent Transformer is required. Notably, this boundary determination must be made without knowledge of bytes that have not yet been generated, necessitating that the patch segmentation mechanism does not rely on forthcoming bytes to partition the sequence. A patching scheme $f_p$ is said to satisfy the criterion of incremental patching if the following holds:
$$f_p(\pmb{x}_{<i}) = f_p(\pmb{x})_{<i}$$
The bpe scheme does not meet the requirements of incremental patching because the way it tokenizes a given prefix can be influenced by the bytes that follow, thereby violating the condition stated above.\footnote{Inserting a dedicated delimiter token to mark patch boundaries can allow bpe to function as an incremental patching scheme, but this modification results in a longer byte sequence.}

\section{{\textsc BLT} Architecture}
\label{section:architecture}
The {\textsc BLT} system is structured around a primary, large-scale autoregressive language model that processes patch-level representations, together with two more compact local models. These local models are responsible for encoding raw byte sequences into patches and for converting patch representations back into bytes, as depicted in~(\autoref{fig:arch_high_level}).

\subsection{Latent Global Transformer Model}

\textit{The Latent Global Transformer} refers to an autoregressive transformer architecture, denoted as $\mathcal{G}$ and composed of $l_{\mathcal{G}}$ layers, that takes as input a sequence of latent patch representations $p_j$ and produces a corresponding sequence of output patch representations $o_j$. In this work, we consistently use the subscript $j$ to refer to patch indices, and $i$ for indices over bytes. The global transformer employs a block-causal attention mask~\citep{dubey2024llama}, constraining attention to patches up to and including the current one within a given document. Since this component dominates both pre-training and inference computational requirements, selectively invoking it allows for dynamic allocation of computational resources, adjusting the processing effort according to the complexity of different parts of the input or output.

\subsection{Local Encoder}
\label{section:local}

\textit{The Local Encoder Model}, symbolized as $\mathcal{E}$, is an efficient transformer-derived architecture containing $l_{\mathcal{E}} << l_{\mathcal{G}}$ layers. Its central function is to translate an input sequence of bytes $b_i$ into informative patch-level representations, $p_j$. Distinct from the canonical transformer design, this model incorporates a cross-attention mechanism subsequent to each transformer block, enabling the pooling of byte-level features into patches~(\autoref{fig:crossattn}).

Initially, input bytes $b_i$ are embedded through a matrix of shape $\mathbb{R}^{256 \times h_{\mathcal{E}}}$, yielding a set of embeddings $x_i$. These base embeddings may be further enriched by hash-embeddings, as detailed in~(\S\ref{sec:hash-emb}).

A stack of alternating transformer and cross-attention modules~(\S\ref{sec:enc-cross-attn}) subsequently refines these embeddings, resulting in patch representations $p_i$ intended for consumption by the global transformer, $\mathcal{G}$.

Transformer components utilize a \textit{local block causal} attention pattern whereby each byte is allowed to focus on up to $w_{\mathcal{E}}$ previous bytes within a sliding window. This window may extend over dynamic patch boundaries but strictly respects document boundaries. Further information regarding the embeddings and specifics of the cross-attention operation is provided in the following subsections.

\subsubsection{Encoder Hash n-gram Embeddings}
\label{sec:ngram-emb}
\label{sec:hash-emb}
One crucial aspect for generating rich and reliable representations at each position $i$ lies in embedding contextual signals from prior bytes. Within {\textsc BLT}, this is accomplished by representing not only the current byte $b_i$ in isolation, but also its inclusion within an n-gram sequence. Specifically, at every step $i$, we begin by assembling the relevant byte-grams

\begin{align}
    g_{i,n}=\{b_{i-n + 1},\ldots, b_i\}
\end{align}

for each byte position $i$ and $n$ from three to eight.\footnote{
We omit byte-grams of size $n$ or more when $i<n$. }

Next, we present hash $n$-gram embeddings, which use a hash function to assign each byte $n$-gram to a corresponding index within a fixed-size embedding table $E_{n}^{hash}$, considering all $n$-gram lengths where $n \in\{3,4,5,6,7,8\}$~\citep{bai2010learning}. The embedding thus obtained is combined with the byte embedding prior to normalization, after which it is fed as input to the local encoder model. To obtain the augmented embedding, we compute

\begin{align}
e_i &= x_i + \sum_{n = 3,...,8}E_{n}^{hash}(\text{Hash}(g_{i,n})) \\
\text{where, Hash}(g_{i,n}) &= \text{RollPolyHash}(g_{i,n}) \% |E_{n}^{hash}|
\end{align}

We scale $e_i$ by dividing it by the total number of $n$-gram sizes incremented by one, utilizing RollPolyHash as specified in Appendix~\ref{appendix:rollpolyhash}. Section~\ref{section:ablations} presents ablation studies examining the influence of $n$-gram hash embeddings, varying both $n$ and the embedding table size, on scaling law patterns under fixed computational costs. Beyond using hashed $n$-gram embeddings, our experiments also encompassed frequency-based $n$-gram embeddings; comprehensive information about these investigations can be found in Appendix~\ref{appendix:ngrams}.

\subsubsection{Encoder Multi-Headed Cross-Attention}
\label{sec:enc-cross-attn}
Our implementation largely mirrors the input cross-attention mechanism from the Perceiver architecture~\citep{jaegle2021perceiver}, though it diverges in that latent vectors are aligned with variable patch representations rather than a fixed set of latent vectors~(\autoref{fig:crossattn}), and these latents only attend to the bytes contained within their respective patch. Specifically, for each patch $p_j$, the cross-attention module constructs a query by first pooling the byte embeddings belonging to $p_j$ and then projecting them via a linear map, $\mathcal{E}_{C} \in \mathbb{R}^{h_{\mathcal{E}} \times (h_{\mathcal{E}} \times U_{\mathcal{E}})}$, with $U_{\mathcal{E}}$ representing the number of heads in the encoder’s cross-attention layer. Formally, let $f_{\text{bytes}}(p_j)$ designate the byte sequence of patch $p_j$; then we define:

\begin{align}
P_{0,j} &= \mathcal{E}_{C}(f_\text{bytes}((p_j)), f~\text{is a pooling operation} \\
P_l &= P_{l-1} + W_o\left(\text{softmax}\left(\frac{QK^T}{\sqrt{d_k}}\right)V\right) \\
\text{where } Q_j &= W_q(P_{l-1,j}), K_i = W_k(h_{l-1,i}), V_i = W_v(h_{l-1,i}) \\
h_l &= \text{Encoder-Transformer-Layer}_l(h_{l-1})
\end{align}

Here, $P \in \mathbb{R}^{n_p \times h_{\mathcal{G}}}$ encodes $n_p$ patch representations, serving as input to the global model. The initial values for $P$ are formed by pooling the byte embeddings $e_i$ belonging to each patch $p_j$. The matrices $W_q$, $W_k$, $W_v$, and $W_o$ are employed as the projection matrices for queries, keys, values, and outputs, respectively, where both keys and values arise from projecting byte-level representations $h_i$ from the preceding layer (for the first layer, $e_i$ are used directly). We adopt a masking mechanism tailored for the patching process: each query $Q_j$ is constrained to attend only to keys and values derived from bytes associated with patch $j$. Given that multi-headed attention is applied to $Q$, $K$, and $V$, and since the patch embedding size $h_{\mathcal{G}}$ typically exceeds $h_{\mathcal{E}}$, during cross-attention the matrix $P_l$ is structured as several attention heads of size $h_{\mathcal{E}}$ before being concatenated back to $h_{\mathcal{G}}$ dimensions. Moreover, pre-LayerNorm normalization is applied to queries, keys, and values in this module, and positional embeddings are omitted within this cross-attention design. A residual connection is also employed to wrap the cross-attention sublayer.

\subsection{Local Decoder} 

The local decoder $\mathcal{D}$, much like its encoder counterpart, utilizes a transformer architecture but remains lightweight by employing significantly fewer layers, such that $l_{\mathcal{D}} << l_{\mathcal{G}}$. Its purpose is to reconstruct raw bytes $y_i$ from an input sequence of global patch representations $o_j$. As it generates each byte, the local decoder bases its predictions on the previously decoded output, drawing upon the hidden states furnished by the local encoder. The decoder structure alternates $l_{\mathcal{D}}$ times between cross-attention modules and standard transformer blocks. In each stack, the cross-attention layer is executed first to map patch-level representations to their corresponding byte-level embeddings, followed by the application of the transformer block directly on this emergent byte sequence.

\subsubsection{Decoder Multi-headed Cross-Attention} 

Within the decoder's cross-attention mechanism, the functions of queries and keys/values are reversed: the byte representations serve as the queries, whereas the patch representations provide the keys and values. The initial input for cross-attention is established as the byte embeddings derived from the final encoder layer, denoted by $h_{l_{\mathcal{E}}}$. For each subsequent decoder layer $l$, the byte representation $d_{l,i}$ is updated according to the following equations:

\begin{align}
D_0 &= h_{l_{\mathcal{E}}} \\
B_l &= D_{l-1} + W_o\left(\text{softmax}\left(\frac{QK^T}{\sqrt{d_k}}\right)V\right), \\
  &\text{where } Q_i = W_q(d_{l-1,i}), K_i = W_k(\mathcal{D}_C(o_j)), V_i = W_v(\mathcal{D}_C(o_j)) \\
  D_l &= \text{Decoder-Transformer-layer}_l(B_l)
\end{align}

Here, as before, $W_k$ and $W_v$ serve as the projection matrices for keys and values, performing linear transformations and employing the split operation $\mathcal{D}_C$ over the patch outputs $o_j$ produced by the global model. The query projection matrix $W_q$ acts on the previous layer’s byte-level representations $d_{l-1}$ in the decoder transformer (or on $h_{l_{\mathcal{E}}}$ for the initial layer). The final output is projected through $W_o$. As a result, the matrix $B$ has dimensionality $\mathbb{R}^{h_{\mathcal{D}} \times n_b}$, with $n_b$ denoting the total number of predicted output bytes. Subsequently, new decoder representations $D_l$ are obtained by passing $B$ through a decoder transformer layer. Consistent with the procedure for local encoder cross-attention, we utilize multi-head attention, apply pre-LayerNorm normalization, omit positional encodings, and employ a residual connection wrapping the cross-attention submodule.

\section{Experimental Setup}
\label{section:experiments}
We perform meticulously controlled experiments to rigorously compare {\textsc BLT} against models based on traditional tokenization, ensuring that {\textsc BLT} does not benefit from potentially extended sequence context lengths.

\subsection{Pre-training Datasets} In all experiments described in this work, each model scale undergoes pre-training on two primary datasets: 1) The Llama 2 dataset~\citep{touvron2023llama}, containing 2 trillion tokens sourced from diverse publicly accessible materials, followed by rigorous cleaning and filtering processes aimed at enhancing data quality; and 2) {\textsc BLT}-1T: a newly curated dataset of 1 trillion tokens, assembled from assorted public datasets and incorporating selected portions of the pre-training data made available by Datacomp-LM~\citep{li2024datacomp}. The Llama 2 dataset underpins our scaling law analyses, leveraging the findings of~\cite{dubey2024llama} to inform optimal token counts and architectural strategies for {\textsc BLT}. In contrast, {\textsc BLT}-1T supports comprehensive pre-training runs, facilitating direct performance comparisons with Llama 3 across downstream evaluation tasks. Importantly, neither dataset draws upon data originating from Meta products or services. Additionally, in establishing baselines for tokenizer-based approaches, we employ the Llama 3 tokenizer configured with a 128K-token vocabulary, as this configuration yielded superior baseline results compared to the Llama 2 tokenizer in our assessments.

\subsection{Entropy Model} For our experiments, the entropy model utilized is a byte-level language model that is trained using the identical data distribution as the full {\textsc BLT} model. By default, the architecture comprises a 100M-parameter transformer consisting of 14 layers, a hidden size of 512, and employs sliding window attention spanning 512 bytes. All other hyperparameter choices align with those adopted for our local and global transformers. We also conducted experiments across various model capacities, receptive fields, and architectural designs, as described in \autoref{section:ablations}. Notably, in cases where the model's receptive field is sufficiently limited, the resulting trained entropy model can be effectively represented via a compact lookup table.

\subsection{Entropy Threshold and Equalizing Context Length} For entropy-based patching models, we determine a patching threshold designed to yield a target average \textit{patch size} over the pretraining dataset. In {\textsc BLT}, as opposed to standard tokenization, the \textit{patch size} is a flexible parameter, which notably affects the context window accessible to the model. To ensure models do not gain an unfair benefit from increased patch sizes by seeing more context, we keep the expected total number of bytes per batch unchanged. Consequently, models configured with larger patch sizes are assigned shorter sequence lengths. Specifically, for Llama 2 data, an 8k byte context is employed, whereas with the {\textsc BLT}-1T corpus, the context is expanded to an average of 16k bytes, all while preserving the batch size at an average of 16M bytes.

Although the overall batch size remains unchanged, dynamic patching approaches produce varying bytes-to-patches ratios across data batches. To enhance efficiency, our approach to {\textsc BLT} training organizes batches by the number of patches, thereby eliminating the need for padding operations in the resource-intensive latent transformer stage. This guarantees a fixed patch count per batch. During the training process, we pad—or, in certain cases, truncate—byte sequences to 12k and 24k bytes for the Llama 2 and {\textsc BLT}-1T datasets, respectively, to prevent excessive memory consumption arising from sequences containing exceptionally large patches.

\subsection{Entropy Model Context}
\label{section:entropy-context}
Our empirical observations indicate that applying entropy patching tends to result in increasingly larger patches, particularly in highly structured tasks such as multiple choice items (refer to \autoref{fig:mmlu_patching} for an example from MMLU), where repetitive patterns are common. This behavior can be attributed to decreased entropy arising from the repetition present within the entropy model context. Consequently, in the extensive {\textsc BLT}-Entropy experiment employing a patch size of 4.5, we mitigate "entropy drift"—which is linked to variations in context length—by reinitializing the entropy context using new lines and enforcing an approximate monotonicity constraint. Notably, this modification is restricted to the calculation of entropy values; we maintain the original methodology when determining the appropriate entropy threshold.

\subsection{FLOPs Estimation} 
\label{section:flops}

Our approach closely adheres to the methodology outlined by Chinchilla~\citep{hoffmann2022training} for calculating transformer flops, accounting for the operations in the feed-forward modules, the qkvo projections within the self-attention mechanism, and the processes involved in computing attention as well as the output projection. However, our analysis departs from theirs in that we regard the input embedding layer as an efficient lookup table rather than a standard dense matrix multiplication, rendering it effectively a 0-flop computation. Consistent with established literature, we also assume the computational cost of the backwards pass is double that of the forwards pass.

To evaluate the flops \textit{per byte} metric for {\textsc BLT} architectures, we aggregate the floating point operations required by the local encoder transformer, the global latent transformer, and the local decoder transformer, in addition to those incurred by the cross attention mechanisms within both the encoder and decoder components:

\begin{align}
    \text{FL}_{\text{{\textsc BLT}}} &= \text{Transf. FL}(h_{\mathcal{G}}, l_{\mathcal{G}}, m=n_{ctx}/n_p,V=0) / n_p\\
   &+ \text{Transf. FL}(h_{\mathcal{E}}, l_{\mathcal{E}}, m=w_{\mathcal{E}}, V=0) \\
   &+ \text{Transf. FL}(h_{\mathcal{D}}, l_{\mathcal{D}}, m=w_{\mathcal{D}}, V=256) \\ 
    &+ \text{Cross Attn. FL}(h_{\mathcal{E}}, l_{\mathcal{E}}, m=n_p, r=n_p/k) \times k/n_p \\ 
   &+ \text{Cross Attn. FL}(h_{\mathcal{D}}, l_{\mathcal{D}}, m=k, r=k/n_p) 
\end{align}

Here, $n_{ctx}$ denotes the sequence length in bytes, $n_p$ indicates the patch size, $r$ represents the ratio of queries to keys/values, and $k$ signifies the proportion of the patch dimension to the byte dimension—that is, the count of local model partitions that are concatenated to construct a global model embedding ($k=2$ in \autoref{fig:crossattn}). The vocabulary size of the output projection is indicated by $V$, which appears solely in the local decoder. The specific value for the attention context length, $m$, depends on whether the module processes a sequence of bytes or patches. Attention flops are subsequently adjusted for each model part based on this distinction. The formulas detailing the calculation of flops for Transformer modules and Cross-Attention are detailed in Appendix~\ref{section:flops-appendix}.

\subsection{Bits-Per-Byte Estimation} Perplexity is meaningful only when a specific tokenizer is used, since it quantifies the model’s uncertainty per token. In order to compare models trained at the byte and token levels, we adopt the Bits-Per-Byte (BPB) metric, as in prior studies~\citep{xue2022byt5, yu2023megabyte, wang2024mambabyte}. BPB serves as a tokenizer-agnostic alternative to perplexity. More formally:

\begin{align}
    \text{BPB}(x)&= \frac{\mathcal{L}_{CE}(\pmb{x})}{\ln(2)\cdot n_{\text{bytes}}}
\end{align}

where the uncertainty over the data $\pmb{x}$ as measured by the sum of the cross-entropy loss is normalized by the total number of bytes in $\pmb{x}$ and a constant. 

\subsection{Transformer Architecture Hyperparameters} In configuring all transformer blocks within {\textsc BLT}—including both the local and global components—we primarily adopt the Llama~3~\citep{dubey2024llama} design. Specifically, our feed-forward networks leverage the SwiGLU activation function~\citep{swiglu}, while rotary positional embeddings (RoPE)~\citep{RoPe} are applied in the self-attention modules with a scaling factor of $\theta=500000$~\citep{xiong2024effective}. For layer normalization throughout, RMSNorm~\citep{RMSNorm} is employed. All self-attention layers that utilize standard attention masks such as \textit{block causal} or \textit{fixed-window block causal} benefit from Flash attention~\citep{dao2022flashattention}; in these cases, the window size for fixed-width attention is set to 512. Flex Attention\footnote{\url{https://pytorch.org/blog/flexattention}} is chosen for our cross-attention layers due to their reliance on dynamically constructed, patch-based masks, resulting in efficient fused kernel implementations and substantial acceleration during training.

\subsection{{\textsc BLT}-Specific Hyperparameters} 

To evaluate the performance of {\textsc BLT} models, we perform experiments in two primary directions: analyzing scaling behaviors and assessing results on downstream tasks. Our experimental setup includes model sizes of 400M, 1B, 2B, 4B, and 8B parameters. The detailed architectural hyperparameters corresponding to these configurations are provided in Appendix~Table~\ref{tab:arch}. 

For initializing the queries in the first cross-attention layer of the local encoder, we employ max-pooling. Across all {\textsc BLT} models, $500,000$ hashes are used, applying a single hash function, and the n-gram sizes span from 3 to 8. All models are trained with a fixed learning rate of $4e-4$. This consistent choice is the outcome of a hyperparameter sweep between $1e-3$ and $1e-4$ conducted at both 400M and 1B scales, which revealed the optimality of the same value for both token-based and {\textsc BLT} models. When analyzing scaling properties on Llama-2 data, we adhere to batch sizes suggested by~\cite{dubey2024llama} or their byte-equivalent where applicable. For optimization, we rely on AdamW~\citep{adamw} with $\beta_1 = 0.9$, $\beta_2 = 0.95$, and $\epsilon=10^{-8}$. The training employs a linear warm-up over 2000 steps, followed by cosine learning rate decay down to zero. Weight decay is set at 0.1, and the gradient norm is clipped globally at 1.0.

\section{Scaling Trends}
\label{section:scaling}

We offer a comprehensive overview of scaling behaviors in byte-level models to guide future scaling efforts for {\textsc BLT} models. This scaling analysis seeks to overcome the shortcomings of prior investigations into byte-level architectures by: (a) comparing patterns within the compute-optimal training regime, (b) training comparable 8B parameter models on substantial datasets (as large as 1T tokens or 4T bytes) and assessing their performance on downstream benchmarks, and (c) evaluating scaling patterns under fixed inference cost conditions. Later in the paper, we turn to a detailed examination of benefits associated specifically with modeling sequences at the byte level.

\subsection{Parameter Matched Compute Optimal Scaling Trends}
\label{section:parameter-matched-scaling}
We conduct experiments with the Llama 2 dataset by training both \textit{compute-optimal} bpe and {\textsc BLT} models, spanning four model scales from 1B to 8B parameters. For each model, we chart the relationship between the amount of training flops expended and its language modeling accuracy on a selected subset of the training mixture. The bpe baselines adhere to the model-parameter-to-data ratio identified as optimal in Llama~3~\citep{dubey2024llama}, an approach specifically intended, per \textit{compute-optimal} theory, to maximize performance for a fixed training compute budget~\citep{hoffmann2022training}, thus serving as a strong reference point. Each bpe model’s counterpart is a {\textsc BLT} model trained on identical data, employing a Latent Transformer that is parameter- and architecture-matched to its bpe baseline.

As shown in~\autoref{fig:scaling-compute-opt} (right), {\textsc BLT} models consistently perform as well as or better than their bpe equivalents, a pattern that persists across increasing model capacities and higher compute (measured in flops). To our knowledge, {\textsc BLT} represents the first instance of a byte-level Transformer matching the scaling behaviors of BPE-based models in compute-optimal settings. This observation thus supports our hypothesis that the ideal ratio between parameters and training compute identified for bpe models similarly applies to {\textsc BLT}, or at the very least, deviates only minimally.

Enhancements in both model architecture and dynamic patching are essential to achieve scaling patterns comparable to bpe. As depicted in ~\autoref{fig:scaling-compute-opt} (left), we evaluate models utilizing space-patching strategies against the performance of Llama 3. Our approximation of SpaceByte \citep{slagle2024spacebyte} employs {\textsc BLT} space-patching, but omits n-gram embeddings and cross-attention mechanisms. While SpaceByte yields better results than Megabyte, its performance does not reach that of Llama 3. On the right panel of \autoref{fig:scaling-compute-opt}, we show the combined effects brought about by architectural modifications and dynamic patching. Models employing {\textsc BLT} demonstrate parity with leading tokenizer-based models, such as Llama 3, when scaling up.

Additionally, we find that the tokenizer selection significantly influences outcomes for tokenizer-based models. Specifically, models trained with the Llama-3 tokenizer demonstrate superior performance relative to those trained on identical data with the Llama-2 tokenizer.

In summary, our {\textsc BLT} framework displays characteristics intermediate between those of Llama 2 and Llama 3 when evaluated with much larger patch sizes. Notably, the average token sizes produced by the bpe tokenizers of Llama 2 and 3 are 3.7 and 4.4 bytes, respectively. By contrast, {\textsc BLT} maintains comparable scaling behaviors with average patch sizes of 6 or even 8 bytes. Since the inference flops decrease as patch size increases, employing an 8-byte patch size results in almost 50\% reduction in inference flops. Furthermore, models that utilize larger patch sizes appear to benefit more from increases in model and data scale. For instance, while {\textsc BLT} with an 8-byte patch size initially performs substantially worse than bpe Llama 2 at a 1B scale, it surpasses bpe's performance at the 7B scale. This trend implies that employing such or even greater patch sizes may be advantageous as models and training compute continue to expand.

\subsection{Beyond Compute Optimal Task Evaluations}
\label{section:evals}
To further investigate scaling dynamics, we extend training of an 8B {\textsc BLT} model past the compute-optimal threshold using the {\textsc BLT}-1T dataset—a larger and higher-quality corpus—and assess its effectiveness across a diverse array of well-established classification and generation benchmarks. For our evaluation, we focus on tasks spanning common sense reasoning, general world knowledge, and code generation:

\paragraph{Classification tasks} comprise ARC-Easy (zero-shot)~\citep{Eval_ARC}, ARC-Challenge (zero-shot)~\citep{Eval_ARC}, HellaSwag (zero-shot)~\citep{Eval_hellaswag}, PIQA (zero-shot)~\citep{Eval_piqa}, and MMLU (five-shot)~\citep{hendrycks2020measuring}. Here, a prompt-based scoring approach is used in which likelihoods over possible choice characters are computed; mean accuracy is then presented.

\paragraph{Coding related generation tasks:} For evaluating code generation abilities, we provide pass@1 metrics on MBPP (three-shot)~\citep{Eval_MBPP} and HumanEval (zero-shot)~\citep{Eval_HumanEval}, specifically measuring LLM-generated Python code quality.

In \autoref{tab:evals}, we present a comparison among three models, each trained on the {\textsc BLT}-1T dataset: a Llama 3 tokenizer-based bpe model\footnote{We select the Llama 3 tokenizer, featuring a 128k vocabulary, due to its superior performance compared to the 32k vocabulary of Llama 2.}, and two {\textsc BLT} model variants. The first variant, {\textsc BLT}-Space, implements a space-patching approach, while the second, {\textsc BLT}-Entropy, adopts an entropy-based patching method. For the entropy-based model, we also enforce an approximate monotonicity constraint and refresh the model’s context with line breaks, as outlined in~\autoref{section:entropy-context}. All three configurations are trained to the same flop budget. Additionally, for {\textsc BLT}-Entropy, we reduce the entropy threshold from 0.6 to 0.1 during inference, a modification that enhances task outcomes, albeit increasing the number of inference steps required.

The {\textsc BLT}-Entropy model achieves superior results compared to the Llama 3 model on 4 out of 7 tasks, despite both models being trained on an equal amount of data in bytes. This enhanced performance can be attributed to two key factors: (1) more efficient utilization of computational resources through dynamic patching, and (2) explicit modeling at the byte level rather than at the token level.

Conversely, {\textsc BLT}-Space lags behind the Llama 3 tokenizer in performance across all but one task; however, due to its larger mean patch size of 6 bytes, it achieves a notable reduction in inference flops. By contrast, models based on bpe and entropy-patching display similar average patch sizes, both around 4.5 bytes on the mixed training dataset. Given an equivalent training budget, the larger patch size enables that model to process 30\% more data relative to the other two methods, potentially causing {\textsc BLT} to diverge further from compute-optimality.

\subsection{Patches Scale Better Than Tokens} In the {\textsc BLT} framework, it is possible to scale up both the model size and the patch size in parallel while holding the compute requirements for training and inference constant, along with the volume of training data. The capability to enlarge patch sizes without bound distinguishes patch-based models from their token-based, fixed-vocabulary counterparts, allowing them to overcome the inherent efficiency constraints described in Section~\ref{section:bpe-patch}. By adopting longer patch sizes, computational costs are reduced, enabling these saved resources to be redirected toward expanding the global latent transformer, since this component is invoked with lower frequency.

In this work, we perform a fixed inference scaling experiment to investigate whether increasing model size while reducing the number of steps over larger patches yields improved performance, compared to using smaller models with more steps. Using base models of 400m and 3.6B parameters paired with the Llama 2 tokenizer, we identify flop-matched counterparts using the Llama 3 tokenizer as well as {\textsc BLT}-Entropy models featuring average patch sizes of 6 and 8 bytes across the training datamix (refer to~\autoref{tab:fixed-inf-params} for detailed model configurations). Notably, the models with patch size 8 are designed with 3 encoder layers instead of just 1. Each configuration is trained under multiple training flop budget settings.

\autoref{fig:fixed_inference_scaling} illustrates that {\textsc BLT} architectures consistently exhibit superior scaling behavior compared to tokenization-based models across both categories of inference compute. While BPE models demonstrate an initial advantage under constrained training budgets, they are soon overtaken by {\textsc BLT} models shortly after surpassing the compute-optimal point. From an applied perspective, allocating additional resources toward pretraining is often advantageous, resulting in models with improved performance at a constant inference cost. The 8B parameter models, such as Llama 3.1, exemplify this phenomenon; these models are trained on datasets that are two orders of magnitude larger than what is considered compute-optimal for their scale.

The point at which {\textsc BLT} surpasses token-based models occurs at a slightly lower multiple of the compute-optimal budget for the models in the larger flop class—shifting from three times to roughly 2.5 times the optimal compute. Additionally, models with patch size 8 in this class display a steeper scaling trend, allowing them to outperform the other models at an earlier point. As noted in Section~\ref{section:parameter-matched-scaling}, increasing the patch size tends to improve {\textsc BLT}'s performance relative to BPE models as the overall model size grows. We partly explain this by the diminishing proportion of total flops consumed by the byte-level Encoder and Decoder components, which do not scale up as rapidly as the Latent Transformer. For instance, when scaling the total parameter count by a factor of 20—from 400M to 8B—the parameter count for {\textsc BLT}’s local components increases by only about twofold. This observation matters because larger patch sizes influence the flops only in the Latent Transformer and leave byte-level module costs unchanged. Indeed, this explains why the parameter size of {\textsc BLT}-Entropy with patch size 8 increases from 1.6x to 1.7x relative to Llama 2 as the model size increases.

In conclusion, our investigation into scaling patch length reveals that the {\textsc BLT} patch-based framework attains superior scaling performance when both the patch size and overall model size are increased together. Notably, these positive scaling effects not only persist but may further intensify as the models scale up.

\section{Byte Modeling Improves Robustness} We further evaluate the robustness of {\textsc BLT} against token-based models that do not have access to explicit byte-level representations. Additionally, we introduce a method for incorporating byte-level processing into existing pretrained token-based models.
\label{sec:robustness}

\subsection{Character-Level Tasks}

Initial efforts to train models at the byte level were largely inspired by their demonstrated resilience to noise introduced at the byte granularity, as well as their sensitivity to the internal structure of tokens—a feature where tokenizer-dependent approaches typically fall short. To empirically assess these aspects, we conduct a series of supplementary experiments utilizing benchmarks specifically designed to test resistance to noisy inputs and sensitivity to character-level information in both English and other languages, encompassing not only alphabetic characters but also numerals and phonetic symbols. The outcomes of these evaluations are summarized in Table~\ref{tab:char_tasks}.

\paragraph{Noisy Data} To evaluate model robustness, we generate corrupted variants of the benchmark classification tasks detailed in Section~\ref{section:evals}, allowing a direct comparison between tokenizer-based models and {\textsc BLT}. Five character-level perturbation methods are utilized to modify the textual inputs: (a) \textit{AntSpeak}: Renders the text fully uppercase with individual characters separated by spaces. (b) \textit{Drop}: Randomly deletes 10\% of characters from the sequence. (c) \textit{RandomCase}: Randomly assigns either uppercase or lowercase to 50\% of the characters each. (d) \textit{Repeat}: Selects 20\% of characters and repeats each up to four times. (e) \textit{UpperCase}: Converts all characters in the text to uppercase form. For each strategy, modifications are independently applied to the prompt, the completion, or both, with each configuration evaluated as a separate task. We aggregate the performance scores by averaging across these scenarios. As shown in Table \ref{tab:char_tasks}, evaluation on the noised HellaSwag~\citep{Eval_hellaswag} dataset demonstrates that {\textsc BLT} consistently surpasses tokenizer-based approaches in robustness, achieving an average improvement of 8 points over a model trained on identical data, and even outperforms the substantially larger Llama 3.1 model.

\paragraph{Phonology - Grapheme-to-Phoneme (G2P)} We evaluate how effectively {\textsc BLT} can convert a word's sequence of graphemes (alphabetic symbols) into its phonemic representation (the corresponding pronunciation symbols). Table \ref{tab:char_tasks} shows the outcomes for the G2P task in a 5-shot framework, using Phonology Bench~\citep{suvarna2024phonologybench}. Our findings indicate that {\textsc BLT} surpasses the performance of the Llama 3 1T tokenizer-based baseline model for this particular task.

\paragraph{CUTE}
To evaluate the ability of models to process information at the character level, we conduct experiments with {\textsc BLT} on the CUTE benchmark~\citep{edman2024cute}. This suite includes multiple tasks grouped into three overarching types: analyzing compositionality, detecting orthographic resemblance, and manipulating sequences. For most models relying on tokenization, CUTE presents substantial difficulties. Although such models typically encode their tokens' orthographic forms, they often fail to effectively leverage these details when manipulating textual inputs. As indicated in Table~\ref{tab:char_tasks}, {\textsc BLT}-Entropy achieves a margin of over 25 points above the BPE Llama 3 models on CUTE. Notably, our approach excels in character-level manipulation, securing scores as high as 99.9\% on spelling-related challenges. These considerable performance gains, despite {\textsc BLT} training on just a sixteenth of the data used for Llama 3.1, underscore the challenges BPE models encounter in acquiring robust character-level representations. Figure \ref{fig:cute_examples} depicts selected cases where the Llama 3 tokenizer-based architecture falls short while our {\textsc BLT} model succeeds. The only areas where BPE-based models outperform are tasks involving word deletion and insertion. While such operations may be more intuitive for token-based methods than for byte-level models, the disparity is minimal, and it may be easier for character-level models to generalize upwards from characters to words, rather than the inverse. Across all tasks, we adhere to the original Huggingface prompts and uniform evaluation criteria. It's plausible that the BPE-based systems could see gains from further prompt optimization.

\paragraph{Low Resource Machine Translation} We assess {\textsc BLT} by conducting translation experiments involving both directions for six widely recognized language families, as well as twenty-one low-resource languages using various scripts, based on the FLORES-101 benchmark~\citep{goyal2022flores}. The performance, measured by SentencePiece BLEU scores, is presented in Table~\ref{tab:flores}.
The findings indicate that {\textsc BLT} achieves superior results compared to a system utilizing the Llama 3 tokenizer, with an overall gain of 2 BLEU points for translations into English and an improvement of 0.5 BLEU points for translations from English. For widely spoken language pairs, the performance of {\textsc BLT} matches or marginally exceeds that of Llama 3. Notably, {\textsc BLT} delivers stronger results than Llama 3 across a variety of low-resource language pairs, highlighting the advantages of byte-level modeling for robustly handling rare and diverse byte sequences.

\subsection{Training {\textsc BLT} from Llama 3}
\label{sec:distillation}
In this approach, we demonstrate how {\textsc BLT} models can capitalize on existing pre-trained models with tokenization, such as Llama 3.1, to facilitate improved and accelerated convergence during training. This is accomplished by initializing the global transformer parameters of {\textsc BLT} with those from a pre-trained Llama 3.1. For the subsequent training, the weights of the global transformer are updated at a learning rate set to one-tenth of that used for the local encoder and decoder models, and we adhere to the optimal training step count established for Llama 3. Comparative results with a baseline {\textsc BLT} are reported in Table~\ref{tab:distill}. The findings indicate that {\textsc BLT} initialized from Llama 3.1 demonstrates markedly superior performance relative to both the Llama 3 and standard {\textsc BLT} baselines, all trained with an equivalent computational budget. Notably, even when measured against the {\textsc BLT}-Entropy variant (refer to Table~\ref{tab:evals}), which was trained with a far larger dataset of 1T tokens, {\textsc BLT} from Llama 3.1 still delivers better results on the MMLU benchmark. This suggests that this initialization strategy can be highly effective for achieving strong results with substantially reduced training flops.

This approach may be interpreted as converting models that rely on tokenizers into those that operate without them, thus adapting a pre-trained LLaMA 3.1 into a {\textsc BLT} framework. For a thorough evaluation, we present results for the original LLaMA 3.1—trained on 15T tokens—in Table \ref{tab:distill}, juxtaposing its performance with the {\textsc BLT} variant derived from it. While the adaptation leads to marginal decreases in performance on MMLU and HumanEval, the impact is more pronounced on several other benchmarks. These findings indicate that additional research is necessary to more effectively exploit the pre-trained model's capabilities, particularly through refining data mixtures and tuning hyperparameters.

\section{Ablations and Discussion}
\label{section:ablations}
Here, we provide an analysis of ablations conducted to support the selection of the {\textsc BLT} architecture, alongside decisions surrounding the patching approach and hyper-parameter configuration for the 8B-parameter {\textsc BLT} model trained on the {\textsc BLT}-1T corpus.

\paragraph{Entropy Model Hyper-parameters} To investigate the impact of different entropy model capacities and context window sizes on scaling dynamics, we experiment with byte-level entropy transformer variants spanning parameter counts from 1 million to 100 million, while also varying context window lengths between 64 and 512 tokens. We display the relationship between bits per byte (bpb) and the number of training flops using scaling curves generated from our $400m$ and $1b$ {\textsc BLT} models trained on Llama-2 data, as seen in \autoref{fig:entropy_ablation}. Our findings show that increasing both the entropy model size and context length tends to improve scaling performance, although these gains diminish as model size exceeds 50 million parameters.

\paragraph{Types of Patching}

We conduct an ablation study on the four patching methods described in Section~\ref{section:patching}, specifically: 1) Strided Patching with strides set to 4 and 6, 2) Whitespace-based Patching, 3) BPE Tokenizer Patching leveraging the Llama 3 tokenizer, and 4) Entropy-based Patching utilizing a compact byte-level language model.

Although dynamic patching shortens the effective sequence length, we regulate sequence length to ensure that each patching strategy maintains a comparable context window. Across all configurations, models are exposed to an equivalent expected number of bytes per sequence during both training and inference, thereby eliminating potential biases related to handling larger contexts. The outcomes of these ablation studies are displayed in Figure~\ref{fig:scaling-compute-opt}. All alternative patching strategies yield better performance than static patching, with space patching closely matching the efficacy of dynamic entropy-based patching.

In \autoref{tab:evals_ablation}, we report evaluation results for {\textsc BLT} models, providing a comparison between models employing tokenizers, those using space patching, and those leveraging entropy-based patching, all of which were trained on the Llama 2 dataset for the empirically optimal number of training steps~\citep{dubey2024llama}. While space patching offers a more straightforward approach that does not require the entropy model to be executed during training, our findings indicate that the performance improvements linked to entropy-based patching seen in scaling studies (see Section~\ref{section:scaling}) are preserved in a variety of downstream benchmark tasks.\footnote{The results for space patching correspond to earlier experiments conducted without cross-attention; nonetheless, the general trends are consistent even with the inclusion of cross-attention.}

\paragraph{Cross-Attention} 
In~\autoref{tab:crossattn_ablations_1b}, we examine the effects of adding cross-attention at different locations within both the encoder and decoder components of {\textsc BLT}. For encoder-side cross-attention, we investigate three types of query initialization: 1) assigning an identical learned embedding to all global states, 2) utilizing a hash embedding derived from the patch bytes, and 3) applying pooling to the encoder's hidden states for the patch bytes at the relevant encoder layer.

Our results indicate that incorporating cross-attention into the \textit{decoder} yields the strongest performance gains. When applied in the encoder, cross-attention provides only modest benefits, which are evident primarily when queries are initialized using pooling. Moreover, the impact of cross-attention is notably pronounced on the Common-Crawl dataset and becomes even more significant as patch size increases.

\paragraph{n-gram Hash Embeddings} 

We experiment with n-gram hash embedding vocabularies set at 0, 100K, 200K, and 400K, and display the corresponding results in Table~\ref{tab:ngram_ablations_1b}. The inclusion of hash embeddings offers improvements across all evaluated domains, with the most pronounced benefits observed for Wikipedia and Github (showing a 0.04 bpb increase, versus a 0.01 bpb difference following 15k steps at the 8B scale). At the 8B model size, reducing the number of hashes from 500K to 300K results in only a marginal performance variation of approximately 0.001 bpb over 15k training steps. This suggests that hash embeddings play a crucial role in enabling {\textsc BLT} to achieve parity with models that employ explicit tokenizers; nevertheless, performance gains become negligible beyond the 300K hash threshold. Moreover, the observed improvements from hashes and cross-attention appear to address different aspects of the modeling task, as their benefits are largely additive across the datasets studied.

\paragraph{Local Model Hyperparameters} Table \ref{tab:local_layers} presents an ablation study on different configurations for the layer count in both the local encoder and decoder. With the incorporation of hash n-gram embeddings, {\textsc BLT} demonstrates optimal performance when the encoder is kept minimal—specifically, a single layer—while the decoder is comparatively more complex.

\section{Related Work}
\label{section:related_work}

\textbf{Character-Level RNNs:} The task of Character Language Modeling has maintained its popularity since the inception of neural approaches~\citep{sutskever2011generating,mikolov2012subword,graves2013generating}, primarily due to the inherent capability of such models to seamlessly handle out-of-vocabulary tokens without the need for explicit back-off strategies. In their work, \cite{kim2016character} present a model that exclusively encodes characters at the input through a combination of convolutional and highway architectures, subsequently integrating this representation into LSTM-driven RNNs. Their approach matches the leading RNN-based language model performances of the era in English, while surpassing them in morphologically rich languages, highlighting a further benefit of character-level language models. Building on this notion, \cite{kenter2018byte} employ byte-level LSTM frameworks for machine comprehension, achieving superior results compared to word-based models, particularly in Turkish and Russian, both of which exhibit rich morphological structure. Likewise, \cite{zhang2015character} explore the use of character-level convolutional networks for classification problems, reporting better performance than their word-based counterparts on certain benchmarks. In addition, \citet{chung2019hierarchical} leverage hierarchical LSTMs enhanced with boundary detectors at various abstraction levels, enabling automatic discovery of latent textual hierarchies to boost character-level language modeling outcomes. Contrasting attention-based approaches, ByteNet~\cite{kalchbrenner2016neural} applies convolutional layers directly to character representations for the task of machine translation.

\textbf{Character-Level Transformers:} The introduction of attention-based transformer architectures~\citep{vaswani2017attention}, along with advances in subword tokenization methods~\citep{sennrich-etal-2016-neural}, substantially advanced neural language modeling and yielded improved results across multiple benchmarks. Nevertheless, both word and subword tokenization strategies embed a particular inductive bias by defining the granularity at which the models process language. In an effort to bridge advances made by transformer models with the positive indications previously reported in character-level language modeling, \citet{al2019character} employed transformers of considerable depth, leveraging auxiliary objectives to train these models. Their approach led to transformer-based character language models that surpassed earlier LSTM-based systems. Yet, their results continued to lag considerably behind those obtained with word-level LLMs. Moreover, findings from GPT-2~\citep{radfordgpt2} demonstrated that, on expansive datasets such as the 1 Billion Word Benchmark, language models operating at the byte level still fell short when compared with their word-level counterparts.

\cite{Choe2019BridgingTG} showed that transformer-based LLMs operating at the byte level can surpass their subword-level counterparts with a similar parameter count; however, these models demand significantly greater computational resources and exhibit substantially longer training times. Likewise, \cite{el2020characterbert} introduced CharFormer, a BERT variant leveraging convolutions over character embeddings to construct word representations. This approach yielded enhanced results within the medical domain, yet incurred considerable computational costs. In another line of work, \cite{clark2022canine} proposed CANINE, an encoder-only model with 150M parameters that processes raw character sequences directly. CANINE employs a deep stack of transformers—functioning akin to our global model—combined with both a local transformer and strided convolution layers to reduce the input sequence length. The model surpassed mBERT, its token-level counterpart, on several downstream multilingual tasks. ByT5~\citep{xue2022byt5} took a different approach by exploring byte-level encoder-decoder architectures that intentionally avoid patching mechanisms. While ByT5 was more robust to noise and demonstrated competitive performance compared to tokenizer-based models using four times less training data, omitting input downsampling necessitated resource-intensive attention calculations over all input bytes, which proved extremely computationally expensive. Because the transition from subword units to byte-level tokens dramatically lengthens input sequences, byte-level models face significant scaling obstacles. The recent advent of the Mamba Architecture~\citep{gu2023mamba}, which maintains a fixed-sized memory across long input contexts, allowed \cite{wang2024mambabyte} to train a byte-level Mamba model—again foregoing patching—that outperformed byte-level transformers under equivalent FLOP conditions at the 350M parameter scale, as measured by bits-per-byte on multiple datasets.

\textbf{Patching-based approaches:} Employing patching techniques can significantly decrease the excessive computational load (measured in flops) associated with byte-level LLMs, potentially preserving model effectiveness. Several studies have reported promising results at modest scales in terms of both model parameters and training corpus size. For example, \cite{nawrot-etal-2022-hierarchical} introduce static patching methods involving downsampling and upsampling, leading to the development of the hourglass transformer. This architecture demonstrates superior performance relative to other byte-level methods at the 150M parameter scale. In follow-up work, \cite{nawrot-etal-2023-efficient} extend these strategies by incorporating dynamic patching methodologies. These include end-to-end learned boundary-predictors, predictors trained with guidance from tokenizers, and an entropy-driven patching approach akin to {\textsc BLT}. Their results reveal that these advancements surpass contemporary vanilla transformers for language modeling on datasets consisting of roughly 400M tokens and models containing 40M parameters. Additionally, \cite{lester2024training} explore sequence compression via arithmetic coding, obtaining compression ratios superior to those yielded by BPE. Through an equal-info windows strategy, they demonstrate improved performance compared to byte-level baselines in language modeling, while still lagging behind methods employing subword units.

Our research takes particular inspiration from MegaByte~\citep{yu2023megabyte}, a decoder-only causal large language model (LLM) employing a fixed, static method for patching and concatenating representations, facilitating the conversion between bytes and patches. On the decoder end, a localized model translates these patches back into bytes. MegaByte has been shown to perform competitively with tokenizer-based models, specifically at the 1B parameter scale, across a corpus comprising approximately 400B bytes. In our experiments, we include comprehensive ablations of MegaByte and observe that static patching falls short compared to contemporary, compute-optimally trained, tokenizer-based approaches when computation is held constant; furthermore, we illustrate how {\textsc BLT} closes this performance gap. Notably, \cite{slagle2024spacebyte} independently identify the same limitations in MegaByte, proposing an enhancement by applying patching to whitespaces and similar bytes, alongside the addition of a local encoder. Their findings show that these adjustments can outperform token-based transformer models within compute-constrained environments on certain data sources, such as Github and arXiv, at the 1B parameter level. In our analysis, we also consider this variant and conclude that to further improve the scalability and effectiveness of byte-level models to a level on par with the leading token-based models—such as Llama 3—additional architectural innovations are still required.

\section{Limitations and Future Work}

In this study, our model training leverages the optimal number of steps as established for Llama~3~\citep{dubey2024llama} to guide our architectural decisions. Notably, these scaling laws are tailored to BPE-level transformers, which suggests that their recommended data-to-parameter ratios may not be directly transferable or ideal for {\textsc BLT}. Deriving scaling laws specifically for {\textsc BLT}, which could yield more advantageous scaling dynamics for our method, is left as a direction for future investigation. Furthermore, the bulk of our experiments utilize models up to 1B parameters; as we explore larger scales, such as 8B parameters and above, the best architectural strategies may evolve, potentially resulting in enhanced outcomes at these higher capacities.

Current transformer libraries and codebases are primarily optimized for the efficiency of tokenizer-based transformer models. Although we conduct experiments that are theoretically matched in terms of flops and leverage some efficient solutions like FlexAttention for handling non-standard transformer layers, our implementations might still fall short of tokenizer-based architectures in actual wall-clock runtime. Additional optimization could therefore further improve our system's performance.

Although {\textsc BLT} currently employs an independently trained entropy model for the patching process, an appealing avenue for further research lies in jointly learning the patching model in an end-to-end manner. In Section \ref{sec:distillation}, we provide preliminary evidence suggesting that it is possible to “byte-ify” tokenizer-based architectures such as Llama 3—which has been trained on over 10T tokens—by leveraging initialization and weight-freezing techniques applied to the global transformer. Progress in this direction may reveal strategies that not only preserve the advantages of bytefying, but also deliver improved performance relative to these tokenizer-based models, eliminating the need to retrain them from the ground up.

\section{Conclusion}
\label{section:conclusion}

This paper introduces the Byte Latent Transformer (\textbf{BLT}), an innovative architecture that challenges the standard reliance on fixed-vocabulary tokenization typically employed in large language models. By adopting a dynamic, learnable strategy for aggregating bytes into patches, {\textsc BLT} enables more adaptive allocation of computational resources in response to varying data complexity. This adaptive mechanism yields notable gains in both computational efficiency and model robustness. Through a comprehensive scaling analysis, we find that {\textsc BLT} attains competitive performance relative to tokenization-based models such as Llama 3, operating effectively up to 8B parameters and 4T bytes of training data. Moreover, it can exchange minimal decreases in evaluation scores for up to 50\% savings in inference computational costs. In addition, {\textsc BLT} introduces a novel scaling opportunity by permitting joint scaling of both model and patch sizes under a constant inference budget—an attractive proposition for real-world compute constraints. Directly modeling raw byte sequences, {\textsc BLT} enhances the model's capacity to manage rare and atypical input patterns, resulting in superior resilience to noisy data and a richer representation of sub-word information. Collectively, these findings establish {\textsc BLT} as a compelling alternative to conventional tokenization-dependent models, offering an efficient, robust, and flexible framework for next-generation language modeling.

\end{document}